\section{Ranks of pure-state reductions}
\label{sec:regional_rank}

\begin{theorem}[Equality of complementary ranks]
  \label{thm:complementary_reduced_ranks}
  \lean{Matrix.rank_partialTraceRight_vecMulVec_eq_rank_partialTraceLeft,
    FiniteProduct.rank_reducedPure_compl, Entropy.rank_regionState_compl,
    Entropy.regionState_eq_reducedPure}
  \leanok
  The two reduced matrices of a bipartite pure vector have equal rank.
  Consequently, for a pure vector on a finite tensor product and a region $R$,
  \begin{align}
    \operatorname{rank}\rho_{R^c}&=\operatorname{rank}\rho_R.
  \end{align}
  No normalization or nonzero-vector assumption is required, and factors of
  dimension zero are allowed.
\end{theorem}
\begin{proof}\leanok
  If $C$ is the coefficient matrix of the vector, the first marginal is
  $CC^\ast$ and the second is the complex conjugate of $C^\ast C$.
  Both Gram matrices have rank $\operatorname{rank}C$. Complex conjugation
  preserves rank, as do transpose and adjoint. Apply this identity in the
  coordinates separating $R$ from its complement.
\end{proof}

\begin{theorem}[Ranks on disjoint unions and between nested regions]
  \label{thm:regional_union_rank}
  \lean{FiniteProduct.rank_reducedMatrix_union_le,
    Entropy.rank_regionState_union_le, Entropy.rank_regionState_sdiff_le}
  \leanok
  Let $\rho$ be positive semidefinite on a finite tensor product. For
  disjoint regions $R,T$,
  \begin{align}
    \operatorname{rank}\rho_{R\cup T}
      &\le\operatorname{rank}\rho_R\operatorname{rank}\rho_T.
  \end{align}
  If the full state is pure and $R\subseteq T$, then
  \begin{align}
    \operatorname{rank}\rho_{T\setminus R}
      &\le\operatorname{rank}\rho_R\operatorname{rank}\rho_T.
  \end{align}
  Neither estimate requires normalization. Empty regions are included.
\end{theorem}
\begin{proof}\leanok
  \uses{thm:rank_product_marginals,thm:complementary_reduced_ranks}
  Reexpress $\rho_{R\cup T}$ on the tensor product of the configuration
  spaces of $R$ and $T$. Its two marginals are $\rho_R$ and $\rho_T$, so the
  bipartite rank inequality gives the first assertion. For the second,
  $(T\setminus R)^c=R\cup T^c$ is a disjoint union. Apply the first inequality
  and equality of the ranks of complementary reductions.
\end{proof}
